\begin{center}
\begin{tabular}{lrrrrr}
\toprule
Сценарий & Потоки & Раунды & Запросы & Успехи & Ожидаемые отказы \\
\midrule
C01 & 2 & 20 & 40 & 20 & 20 \\
C02 & 10 & 20 & 200 & 20 & 180 \\
C03 & 2 & 20 & 40 & 20 & 20 \\
\midrule
Всего & & 60 & 280 & 60 & 220 \\
\bottomrule
\end{tabular}
\end{center}

\begin{itemize}
\item \textbf{C01. Редактирование и удаление одного фильма.}
Два администратора получили одинаковую версию объекта. В 9 раундах сохранилось редактирование, в 11 --- удаление. При победе редактирования проверены слоган и увеличение версии на единицу. При победе удаления повторное чтение подтвердило отсутствие фильма.
\item \textbf{C02. Создание фильма с общим ключом.}
Десять пользователей одновременно отправили одинаковые название и идентификатор режиссёра. В каждом раунде сохранился один фильм с заданным ключом. Остальные запросы получили HTTP 409 с кодом \texttt{DUPLICATE\_MOVIE}.
\item \textbf{C03. Импорт файлов с общей второй записью.}
Каждый файл содержит свой первый фильм и общий второй. Все 20 отклонённых файлов получили конфликт в \mbox{\texttt{movies[1].director}}. В каждом раунде сохранились ровно два фильма из файла-победителя. Снимки PostgreSQL подтвердили прирост на 40 фильмов, 40 координат, 20 людей и 20 мест. История сохранила успешные и отклонённые операции с соответствующими счётчиками.
\end{itemize}

Все три сценария прошли проверки HTTP-ответов и конечного состояния. Неожиданных отказов, ошибок JTL, повторяющихся ключей и остатков отклонённых импортов не обнаружено.
